\section{Type Inference}
Type inference proceeds in two steps: 

\begin{enumerate}
    \item \textit{Constraint generation}. Traverses the program once to generate constraints over the modes of float typed expressions.
    \item \textit{Constraint solving}. Solves the constraints to produce a fully annotated program.
\end{enumerate}

\subsection{Constraint Generation}
Let the modes in a float type $\Float{m}$ be:
\begin{align*}
    m := E \mid G \mid \alpha
\end{align*}
where $\alpha$ is a mode metavariable. In particular, $\alpha$ is an unknown mode introduced by the inference algorithm. Each fresh $\alpha$ stands for a value in $\{\E, \G\}$ that will be fixed later by solving constraints.

\paragraph{Floats and arithmetic}
\begin{mathpar}
  \inferrule[\textsc{FloatLit}]{\;}{\Gamma \vdash c : \Float{\alpha} \mid \emptyset}
  \and
  \inferrule[\textsc{Add}]
    {\Gamma \vdash e_1 : \Float{m_1} \mid C_1 \\ \Gamma \vdash e_2 : \Float{m_2} \mid C_2}
    {\Gamma \vdash e_1 + e_2 : \Float{\alpha} \mid C_1, C_2, m_1 \preceq \alpha, m_2 \preceq \alpha}
  \and
  \inferrule[\textsc{Mul-G}]
    {\Gamma \vdash e_1 : \Float{m_1} \mid C_1 \\ \Gamma \vdash e_2 : \Float{m_2} \mid C_2}
    {\Gamma \vdash e_1 \times e_2 : \Float{\G}  \mid C_1, C_2, m_1 \preceq \G, m_2 \preceq \G}
  \and
  \inferrule[\textsc{Mul-ConstL}]
    {\Gamma \vdash c : \Float{m_1} \mid C_1 \\ \Gamma \vdash e : \Float{m_2} \mid C_2}
    {\Gamma \vdash c \times e : \Float{\alpha} \mid C_1, C_2, m_1 \preceq \alpha, m_2 \preceq \alpha}
  \and
  \inferrule[\textsc{Mul-ConstR}]
    {\Gamma \vdash e : \Float{m_1} \mid C_1 \\ \Gamma \vdash c : \Float{m_2} \mid C_2}
    {\Gamma \vdash e \times c : \Float{\alpha} \mid C_1, C_2, m_1 \preceq \alpha, m_2 \preceq \alpha}
  \and
  \inferrule[\textsc{Cmp}]
    {\Gamma \vdash e_1 : \Float{m_1} \mid C_1 \\ \Gamma \vdash e_2 : \Float{m_2} \mid C_2}
    {\Gamma \vdash e_1 < e_2 : \Bool : C_1, C_2, m_1 \preceq \G, m_2 \preceq \G}
\end{mathpar}

\paragraph{Distributions}
\begin{mathpar}
  \inferrule[\textsc{Uniform}]
    {\Gamma \vdash e_1 : \Float{m_1} \mid C_1 \\ \Gamma \vdash e_2 : \Float{m_2} \mid C_2}
    {\Gamma \vdash \uniform(e_1,e_2) : \Float{\alpha} \mid C_1, C_2, m_1 \preceq \alpha, m_2 \preceq \alpha}
  \and
  \inferrule[\textsc{Gauss}]
    {\Gamma \vdash e_1 : \Float{m_1} \mid C_1 \\ \Gamma \vdash e_2 : \Float{m_2} \mid C_2}
    {\Gamma \vdash \gaussian(e_1,e_2) : \Float{\alpha} \mid C_1, C_2, m_1 \preceq \alpha, m_2 \preceq \G}
  \and
  \inferrule[\textsc{Poisson}]
    {\Gamma \vdash e_1 : \Float{m_1} \mid C_1}
    {\Gamma \vdash \poisson(e_1) : \Float{\alpha} \mid C_1, m_1 \preceq \alpha}
  \and
  \inferrule[\textsc{Exponential}]
    {\Gamma \vdash e_1 : \Float{m_1} \mid C_1}
    {\Gamma \vdash \exponential(e_1) : \Float{\alpha} \mid C_1, m_1 \preceq \G}
  \and
  \inferrule[\textsc{Beta}]
    {\Gamma \vdash e_1 : \Float{m_1} \mid C_1 \\ \Gamma \vdash e_2 : \Float{m_2} \mid C_2}
    {\Gamma \vdash \betafn(e_1,e_2) : \Float{\alpha} \mid C_1, C_2, m_1 \preceq \G, m_2 \preceq \G}
  \and
  \inferrule[\textsc{Gamma}]
    {\Gamma \vdash e_1 : \Float{m_1} \mid C_1 \\ \Gamma \vdash e_2 : \Float{m_2} \mid C_2}
    {\Gamma \vdash \gammafn(e_1,e_2) : \Float{\alpha} \mid C_1, C_2, m_1 \preceq \alpha, m_2 \preceq \G}
\end{mathpar}

\paragraph{Control}
\begin{mathpar}
  \inferrule[\textsc{If}]
    {\Gamma \vdash e_1 : \Bool \mid C_1 \\ \Gamma \vdash e_2 : \tau_2 \mid C_2 \\ \Gamma \vdash e_3 : \tau_3 \mid C_3}
    {\Gamma \vdash \ifkw\;e\;\thenkw\;e_2\;\elsekw\;e_3 : \tau_4 \mid C_1, C_2, C_3, \tau_2 <: \tau_4, \tau_3 <: \tau_4}
  \and
  \inferrule[\textsc{Let}]
    {\Gamma \vdash e_1 : \tau_1 \mid C_1 \\ \Gamma, x{:}\tau_1 \vdash e_2 : \tau_2 \mid C_2}
    {\Gamma \vdash \letkw\;x=e_1\;\inkw\;e_2 : \tau_2 \mid C_1, C_2}
\end{mathpar}

\begin{example}
    Consider the following program:
    \[
        \letkw\ x = \uniform(0,1)\ \inkw\ \betafn(x, 2)
    \]

    We annotate it with types:
    \[
    \begin{aligned}
      &(\letkw\ x = \uniform(0 : \Float{\alpha_1},\, 1 : \Float{\alpha_2}) : \Float{\alpha_3} \\
      &\quad \inkw\ \betafn(x : \Float{\alpha_3},\, 2 : \Float{\alpha_4}) : \Float{\alpha_5}) : \Float{\alpha_5}.
    \end{aligned}
    \]

    The following constraints are generated:
    \begin{align*}
      \alpha_1 \preceq \alpha_3,\ \alpha_2 \preceq \alpha_3,\ \alpha_3 \preceq \G,\ \alpha_4 \preceq \G.
    \end{align*}

\end{example}

\begin{example}
    Consider the following program:
    \[
    \ifkw\ \flip(0.5)\ \thenkw\ \gammafn(\uniform(0,1), 1)\ \elsekw\ 2 * \betafn(1,4)
    \]
    
    We annotate it with types:
    \[
    \begin{aligned}
    &(\ifkw\ \flip(0.5) \\
    &\quad \thenkw\ \gammafn\!\big(
      \uniform(0 : \Float{\alpha_1},\, 1 : \Float{\alpha_2}) : \Float{\alpha_3},\,
      1 : \Float{\alpha_4}
    \big) : \Float{\alpha_5} \\
    &\quad \elsekw\ \big(
      2 : \Float{\alpha_6}
      *\,
      \betafn(1 : \Float{\alpha_7},\, 4 : \Float{\alpha_8}) : \Float{\alpha_9}
    \big) : \Float{\alpha_{10}} \\
    &)\ : \Float{\alpha_{11}}.
    \end{aligned}
    \]
    
  The following constraints are generated:
  \begin{align*}
      \alpha_1 \preceq \alpha_3 &, \alpha_2 \preceq \alpha_3 \\
      \alpha_3 \preceq \alpha_5 &, \alpha_4 \preceq \G \\
      \alpha_6 \preceq \alpha_{10} &, \alpha_9 \preceq \alpha_{10} \\
      \alpha_7 \preceq \G &, \alpha_8 \preceq \G \\
      \Float{\alpha_5} <: \Float{\alpha_{11}} &, \Float{\alpha_{10}} <: \Float{\alpha_{11}}
  \end{align*}
\end{example}

\begin{example}
    With lists.
\end{example}

\subsection{Constraint Solving}
Constraint solving then finds a substitution mapping each metavariable $\alpha$ to $E$ or $G$. We choose the greatest solution with respect to $\G \preceq \E$, i.e. choose $\E$ whenever the constraints do not force $\G$. At the end of the constraint solving phase, set any remaining $\alpha$ to $\E$.

\begin{example}
  For the generated constraints above, the greatest solution is:
  \begin{align*}
    \theta = [&\alpha_1 \mapsto \G,\ \alpha_2 \mapsto \G,\ \alpha_3 \mapsto \G,\ \alpha_4 \mapsto \G,\ \alpha_5 \mapsto \E ].
  \end{align*}
  In particular, the program is assigned type $\Float{\E}$.
\end{example}

\begin{example}
  For the generated constraints above, the greatest solution is:
  \begin{align*}
    \theta = [&\alpha_1 \mapsto \E,\ \alpha_2 \mapsto \E,\ \alpha_3 \mapsto \E,\ \alpha_4 \mapsto \G,\ \alpha_5 \mapsto \E, \\
              &\alpha_6 \mapsto \E,\ \alpha_7 \mapsto \G,\ \alpha_8 \mapsto \G,\ \alpha_9 \mapsto \E,\ \alpha_{10} \mapsto \E,\ \alpha_{11} \mapsto \E ].
  \end{align*}
  In particular, the program is assigned type $\Float{\E}$.
\end{example}
